\section{预测方法与适用性讨论}
分别对存储0/1的$\Delta Q$和$\Delta M$建立四个岭回归器，特征为
\begin{equation}\mathbf x=(Q_{\rm dep},Q_{\rm dep}/L,Q_{\rm dep}\sin\theta).\end{equation}
八个物理工况留一交叉验证中，每组四次重复及配对位保持同组，标准化仅使用该折训练数据。
候选正则化系数为0、0.01、0.1、1、10，最终参数在测试TCAD首次提交前冻结。
模型是给定Geant4沉积与基线后的条件TCAD响应代理，不代替任意粒子输运。

全量测试扩大了已评价的尺寸范围，却没有校准真实工艺或增加能量、位置、材料栈覆盖。
相同简化模型可能具有共同系统误差，这些误差不会因为训练或测试残差小而消失。
已报告误差须与基准和最差误差一起阅读，不能仅择优列出某个位值或粒子。
低维特征关联不能唯一证明微观收集机制。

存储状态由理想钳位和准静态初态建立，未验证真实写入能力；二维1T1C、固定集总电容和
总量等效直径迹不代表真实立体存储器及三维敏感体积。所选电磁option4的质子结果不包含已验证的
核反应次级粒子效应。4/6 nm结果只是固定硅入口事件下TCAD敏感性，不是完整材料栈输运重算。

连续门槛、判据分类和持久位翻转是不同对象。采样时刻的临界与跨界不能自动解释为商业SEU，
每层N=4也不能支持高精度稀有失效概率。数值细化通过的响应和界外未细化的响应均在完整状态表中区分。
确定性工程重试保留原失败及未知退出码，只恢复原逻辑槽。

AI协助脚本、核验和文稿，固定程序执行数值过程。本研究未设置人工操作对照，因此没有AI加速率结论。
学生与教师活动不是本文软件证据的推论，课程实际审核和答辩保留待完成。
